\appsection{Controlled Vocabulary}

The monograph's named objects, one line each, with the defining location. Terms are grouped by the layer that introduces them; the six \textbf{tracked coordinates} are marked ◆.

{\def\LTcaptype{none} % do not increment counter
\begin{longtable}[]{@{}
  >{\raggedright\arraybackslash}p{(\linewidth - 4\tabcolsep) * \real{0.2918}}
  >{\raggedright\arraybackslash}p{(\linewidth - 4\tabcolsep) * \real{0.2838}}
  >{\raggedright\arraybackslash}p{(\linewidth - 4\tabcolsep) * \real{0.4244}}@{}}
\toprule\noalign{}
\begin{minipage}[b]{\linewidth}\raggedright
term
\end{minipage} & \begin{minipage}[b]{\linewidth}\raggedright
defined
\end{minipage} & \begin{minipage}[b]{\linewidth}\raggedright
one-line meaning
\end{minipage} \\
\midrule\noalign{}
\endhead
\bottomrule\noalign{}
\endlastfoot
contrast domain \(D\) & \hyperlink{sec:2}{§2}, \hyperlink{stmt:4.1}{Def~4.1} & the space of alternatives relative to which information is assessed \\
view; kernel \(\ker(v)\) & \hyperlink{stmt:4.2}{Def~4.2} & a presentation-producing map on \(D\); the partition of candidates it cannot separate \\
registration; canonical registration & \hyperlink{stmt:4.6}{Def~4.6} & the record actually kept of a view's output; canonical = lossless \\
registered content \(\sigma(S)\) & \hyperlink{stmt:5.2}{Def~5.2} & \(\bigvee_{v \in S} \ker(v)\): the distinctions the views jointly draw \\
answerability & \hyperlink{stmt:5.6}{Prop~5.6} & a question is answerable iff its kernel factors below the content \\
achievable quotients \(\mathrm{Ach}(V)\) & \hyperlink{stmt:6.2}{Thm~6.2}(3) & the contents realizable by view sets; the algebra's label lattice \\
piece; combination; focusing & \hyperlink{sec:3.4}{§3.4}, \hyperlink{stmt:6.11}{Prop~6.11} & labeled saturated set \((E, \pi)\); intersection-with-label-join; saturation-with-relabel \\
admissibility structure \(K\); ceiling \(\gamma_K\) & \hyperlink{stmt:14.3}{Def~14.3} & the decoder's constraint: admissible readings and their content cap \\
separable vs joint content & \hyperlink{stmt:15.3}{Def~15.3} & what per-view admissible reading recovers vs joint admissible reading \\
registration anomaly \(\Delta_K(T)\) & \hyperlink{stmt:15.3}{Def~15.3} & the interval from separable to joint admissible content---\hspace{0pt}the layer's obstruction \\
◆ compositionality; ◆ codomain & \hyperlink{sec:24}{§24} table & whether content over unions decomposes; whether content lives in a lattice \\
◆ representability & Rem 24.6 & whether content is a principal down-set (an element) or a general lower set; not necessarily a directed ideal \\
stratum: \(\sigma^G\) / \(\sigma^0\) / \(\sigma^{=}\) / \(\widehat\sigma\) & Defs 21.4--21.6, Prop 22.3′ & graph-valued zero-error content; its partition shadow; likelihood-equality content; the Blackwell class \\
coupling; coupling fiber & Rem 21.3, \hyperlink{stmt:39.1}{Def~39.1} & the joint behavior marginals do not determine; the space of its possibilities \\
accumulation & \hyperlink{stmt:23.3}{Prop~23.3}, \hyperlink{stmt:84.1}{Thm~84.1} & independent looks compound---\hspace{0pt}the failure of idempotent combination \\
presentation / context / verdict; policy & \hyperlink{sec:19}{§19}, \hyperlink{stmt:27.6}{Def~27.6} & the anchored record's three atoms; the selection resolving ambiguous attribution \\
leak & \hyperlink{stmt:28.2}{Prop~28.2}, \hyperlink{stmt:45.2}{Prop~45.2} & verdict distinctions exceeding the presentation's---\hspace{0pt}content carried by attribution \\
blind / targeted pattern & \hyperlink{stmt:29.1}{Def~29.1} & corruption indifferent to the candidate / conditioned on it \\
loyalty anomaly; sterility & \hyperlink{stmt:30.2}{Prop~30.2}, \hyperlink{stmt:30.3}{Def~30.3} & erring policies can outinform loyal ones; contexts that cannot leak \\
◆ fidelity & \hyperlink{sec:32.2}{§32.2} & the record-vs-world axis: provenance, certification, forgery \\
cover; defect; descent & §§35--37 & a decomposition of a view set; the interval its data fail to glue by; the gluing discipline \\
reflection; witness & \hyperlink{stmt:37.7}{Thm~37.7} & component-level intervals from graph reflection and witness loss; missing edges alone need not make either interval nonzero \\
◆ descent (coordinate) & \hyperlink{sec:42}{§42} & which layer's data glue: readings/content (cosheaf side) vs distributions (sheaf side) \\
address & Part VII & the mechanism-and-side location of an information quantum (witness / reflection / coupling) \\
ledger & \hyperlink{sec:52}{§52} & the data held fixed while alternatives range---\hspace{0pt}fine (candidatewise) vs coarse (question-relative) \\
evaluation; rung & \hyperlink{stmt:53.1}{Def~53.1}, \hyperlink{stmt:68.1}{Def~68.1}, Rem 40.3 & the anomaly measured---\hspace{0pt}in bits (\(\nu_{\mu,q}\)) or deficiency (\(\nu_\delta\)); the coefficient level (exact / Shannon / enriched) \\
lever & \hyperlink{stmt:46.3}{Thm~46.3}, \hyperlink{sec:63}{§63} & the policy, as the one interpreter-set parameter that reaches the obstruction classes \\
half-gap & \hyperlink{stmt:69.2}{Prop~69.2} & nondegenerate exact intervals have value \(\ge \tfrac12\); only a strictly positive smaller value certifies grading \\
adaptivity premium; causal gap; commitment & Thms 76.2, 77.1, \hyperlink{stmt:78.2}{Def~78.2} & history's value to the policy; futurity's cost to the adversary; write-ahead as internalized provenance \\
domain family; vacuous extension & Defs 83.1, 83.4 & related contrast domains under coarsening/restriction/products; pullback of pieces \\
◆ algebra (coordinate) & \hyperlink{stmt:84.1}{Thm~84.1} & generalized set-algebra laws and independent-replication idempotency; a full graded valuation algebra remains unproved \\
\end{longtable}
}

\textbf{Symbols with more than one meaning.} A few letters are reused across parts; the meaning is always fixed by the local context, as listed here.

{\def\LTcaptype{none} % do not increment counter
\begin{longtable}[]{@{}
  >{\raggedright\arraybackslash}p{(\linewidth - 2\tabcolsep) * \real{0.1351}}
  >{\raggedright\arraybackslash}p{(\linewidth - 2\tabcolsep) * \real{0.8649}}@{}}
\toprule\noalign{}
\begin{minipage}[b]{\linewidth}\raggedright
symbol
\end{minipage} & \begin{minipage}[b]{\linewidth}\raggedright
meanings, with the parts where they occur
\end{minipage} \\
\midrule\noalign{}
\endhead
\bottomrule\noalign{}
\endlastfoot
\(S\), \(T\) & finite view sets (throughout); in Part IV, \(S\) is also the standard record (presentation and verdict), and view sets there are written \(W\) \\
\(\tau\) & the upper adjoint of the content map (Definition 6.1, \hyperlink{stmt:15.1}{Theorem~15.1}); the verdict view \(\tau_v\) of an anchored view (Parts VI, VIII, X); the separation \(\tau\) of a dichotomy in total variation (Proposition 71.1) \\
\(N\) & the forgetful record (Part IV); the presentation view \(N_v\) (Part VI); a binary symmetric channel (Proposition 23.3, Appendix B) \\
\(\delta\) & Le Cam deficiency \(\delta(E,F)\) (Parts V, IX; Appendix D); point masses \(\delta_y\); a cover defect \(\delta_K(\mathcal U;T)\) (§37); the erasure probability in \hyperlink{stmt:61.1}{Theorem~61.1} and \hyperlink{stmt:71.2}{Proposition~71.2}; the minimum total-variation gap in the proof of \hyperlink{stmt:22.2}{Theorem~22.2} \\
\(\Delta\) & the registration anomaly \(\Delta_K(T)\) (Definition 15.3) and its components (Theorem 37.7); Le Cam distance \(\Delta(E,F)\) (§40); the simplex \(\Delta^k\) (Theorem 38.4); the BROJA feasible set \(\Delta_P\) (Proposition 55.4) \\
\(d\) & the label map \(d(E,\pi)=\pi\) (Proposition 6.11); a distractor origin (Parts IV, VI, VIII); the \(L_1\) separation of two class averages (Lemma 62.1) \\
\(K\) & an admissibility structure (Parts II, V, VI); a regular conditional kernel (Appendix A) \\
\(m\) & the minimal sufficient statistic (Theorem 22.6); the number of verdict symbols (Theorem 63.2); fiber minima \(m_D\), \(m_A\) (Appendix D.4) \\
\(\varepsilon\) & a corruption rate (Parts IV, V, VIII--X); an error threshold in \(\varepsilon\)-answerability (Definition 22.0) \\
\(\iota\) & the irreparable attribution rate (Definition 27.3) \\
\end{longtable}
}

\textbf{One overloaded word, disambiguated.} ``Kernel'' is used in exactly two senses: (i) \(\ker(v)\), the kernel \emph{of a map}---\hspace{0pt}always with its argument; (ii) \textbf{the kernel}, Part I's exact deterministic core---\hspace{0pt}the maximal-structure fiber of the collapse theorem, called the \emph{exact core} where prose risks ambiguity. ``Markov kernel'' also occurs for a stochastic channel and has a different, explicitly qualified meaning.
